\chapter[Designing BSS Station Networks]{Designing \ac{bss} Station Networks}
\label{ch:station_placement}

The spatial configuration of docking stations is one of the most critical decisions in the planning and operation of docked \acp{bss}. The location, density, and distribution of stations fundamentally shape system accessibility, determine potential demand, and influence both operational efficiency and user experience.

Given the complexity and multi-faceted nature of station placement decisions, researchers and practitioners have developed and used a diverse range of approaches to address this challenge. These methods vary in their underlying assumptions, data requirements, optimization techniques and objectives, and computational demands. The choice of method is shaped by factors such as local context, data availability, and computational resources, but ultimately hinges on the objectives guiding the planning process. Because of this, station placement is not merely a technical optimization task, but a strategic decision with direct implications for system performance and user experience.

By synthesizing the station placement body of work, this chapter establishes the conceptual and methodological foundations for the station planning framework developed in Part~\ref{pt:data-driven-station-planning} of this thesis. The following sections, which introduce the station placement optimization problem and provide a comprehensive literature review of existing approaches, are adapted from Grau-Escolano et al.~\cite{GrauEscolano2025}, a manuscript submitted to the \textit{Journal of Urban Mobility} and currently under review.

\section{Introduction}

The strategic placement of stations is a decisive factor of operational success, influencing ridership levels~\cite{NACTO2016, Kabra2018} and cost efficiency of the system~\cite{Fu2022}. Historically, \ac{bss} station planning has relied heavily on heuristics based on expert judgment~\cite{Ferrando2007, Yanocha2018, CROW2006}. Planners commonly identified broad zones of high population density or proximity to major trip generators, such as \ac{pt} stations, employment districts, or shopping areas, and then refined locations iteratively based on budget and qualitative assessments of suitability. Similarly, fleet sizes were set according to rough demand estimates or available budgets. While this approach laid the groundwork for many early systems, it often relied more on intuition and practical experience than on formal analytical methods, and was primarily oriented toward utilitarian needs. Consequently, systematically evaluating competing goals and adapting to changing urban conditions was challenging.

The proliferation of high-resolution spatial data and advanced computational methods has enabled a shift toward data-driven decision-making in \ac{bss} planning. Various methods, including location–allocation models~\cite{Frade2015, Mete2018}, \ac{mcdm}~\cite{Kabak2018, Bahadori2022}, and exact and metaheuristic optimization techniques~\cite{Cintrano2018, Qian2022, Caggiani2020} have been applied to optimize station locations. Yet, regardless of the methodology used, the design and implementation of most docked systems still follow an iterative workflow in which the first and most critical step is the clear definition of planning objectives~\cite{Duran-Rodas2021, Ferrando2007, Yanocha2018}. These may include expanding active mobility~\cite{tfgmActiveTravel2025}, improving first- and last-kilometer connectivity~\cite{MTC2013BayAreaBSS}, reducing \ac{ghg} emissions~\cite{ctcn_velib}, or promoting social inclusion~\cite{tembiciBogota2023}. Such objectives can differ across regions within the same system and often evolve over time; for instance, the expansion of an existing \ac{bss} may prioritize goals that are different from those guiding its original implementation.

Despite advances, existing data-driven methods exhibit limitations: they are often applied to pursue a single dominant planning objective, making it difficult to explicitly model trade-offs among genuinely competing goals, such as maximizing demand and improving social inclusion. For instance, while \ac{mcdm} approaches can incorporate many data sources, they are typically used to optimize a single dominant objective, often corresponding to maximizing coverage~\cite{Fazio2021} or demand~\cite{Kabak2018, Bahadori2022}. Even approaches explicitly designed to handle multiple planning objectives usually limit the analysis to a small, predefined set of goals, most often only two, such as demand and social equity~\cite{Conrow2018, Qian2022, Caggiani2020, Duran-Rodas2021} or social equity and first–last-mile connectivity~\cite{Fan2024}. Moreover, in one identified case~\cite{Nikiforiadis2021} three planning objectives were considered, namely demand, social equity, and cost. Relying on such predefined sets of planning objectives constrains the ability to explore a wider range of scenarios without using multiple models or redesigning the optimization framework.

In addition, the spatial resolution of candidate station locations remains a persistent challenge. Many studies aggregate these into coarse administrative zones or regular grids to reduce computational complexity~\cite{Frade2015, Fazio2021, Veillette2018, Nikiforiadis2021}. Although this simplification can ease model formulation and reduce the risk of overfitting when modelling, it can have its drawbacks. First, it prevents planners from working directly with the street layout, which is the level of detail needed for deciding the final station locations. Second, it can hide important local factors, such as the exact location of trip generators (e.g. workplaces, schools, or public facilities) or nearby cycling infrastructure.

Beyond the spatial resolution of candidate locations, how the transport network is represented becomes equally important, as it shapes the estimation of how users travel between stations and, ultimately, how they use the system. Few studies have gone beyond simply mapping potential sites onto a street-level cycling network to actually integrating the network as a graph into the optimization process, where station placement is evaluated based on a more realistic travel paths and connectivity~\cite{Cintrano2018, Xin2013}. In these cases, the network is represented as an undirected graph, meaning that all links are treated as bidirectional regardless of actual circulation rules. This simplification is already a step toward more realistic modelling compared to not acknowledging the distance between stations or considering purely Euclidean distances, but it can still misrepresent actual cycling conditions. In practice, bikes should follow directed networks that respect one-way links, turn restrictions, and other routing constraints. Incorporating these constraints would enhance realism of how users use the bikes and the evaluation of key metrics, such as station-to-station travel distance.

Furthermore, slope has been found to be a key determinant of cycling behaviour and effort in the literature~\cite{Grau-Escolano2024, MateoBabiano2016}, with different effects on \acp{m-b}and \acp{e-b}~\cite{Grau-Escolano2024}. Despite this, most \ac{bss} placement studies commonly omit topography when optimizing locations~\cite{Fazio2021, Kabak2018, Ebrahimi2022, Tera2023}. This limitation is particularly relevant in cities with pronounced elevation changes, where neglecting slope can distort estimates of station accessibility and user suitability. In this regard, one exception is the work of García-Palomares et al.~\cite{GarciaPalomares2012}, who incorporated slope as a uniform penalty on effective travel time. While this is a relevant step in incorporating topography effects, their approach did not distinguish between uphill and downhill gradients, thereby overlooking the facilitating effect of moderate declines and the safety-related penalties of steep descents. By neglecting these effects, the accessibility and suitability of certain locations can be misrepresented.

Finally, the ability to expand and adapt \ac{bss} over time is essential for ensuring long-term system efficiency and alignment with changing mobility patterns. However, many \ac{bss} planning methods are designed for single-phase deployments, in which the entire network is planned at once for the initial roll-out, and thus lack mechanisms to support incremental expansion or reconfiguration while considering existing \ac{bss} stations~\cite{Kabak2018, Bahadori2022, Duran-Rodas2021}. Although these models can be applied to extend an existing \ac{bss}, they are typically not designed to consider how existing stations are positioned relative to each other and their urban context. As a result, optimizing new stations as part of a cohesive system is often challenging, and they risk being treated as isolated additions disconnected from the existing network.

Taken together, these insights highlight the opportunity to develop planning tools that integrate multiple objectives, fine-grained spatial detail, and realistic modeling of the operational environment. To address these gaps, Part~\ref{pt:data-driven-station-planning} introduces a flexible, data-driven framework for optimizing the placement of \ac{bss} stations. Its main contributions are as follows: 

\begin{enumerate}
    \item  \textbf{A general \ac{mcdm} scoring framework:} a modular scoring approach that quantifies the suitability of each candidate station location by aggregating spatial attributes within adjustable buffers that capture the surrounding conditions of potential sites. 
    
    Unlike models limited to a predefined set of planning objectives, this framework is designed to allow planners to integrate any spatial factor based on the specific planning objectives of the \ac{bss}, and to freely choose their relative weights to reflect any combination of priorities. This flexibility supports the creation of alternative planning scenarios, enabling the design of station networks aligned with multiple objectives without requiring separate models.

    \item \textbf{Integrated network metrics on directed graphs:} The framework incorporates network-level measures of station dispersion and accessibility, enabling planners to balance trade-offs between station clustering and system accessibility explicitly. These metrics are computed on a directed network that respects one-way streets, turn restrictions, and other routing constraints, and consider all intersections of the bicycle network as potential station locations. This provides a realistic basis for evaluating and comparing configurations.

    \item \textbf{Topography-adjusted distance modeling:} to improve the accuracy of proximity and accessibility assessments, the model optionally integrates elevation-adjusted distance calculations that reflect slope-induced variations in cycling effort. This feature is particularly important in cities with significant topographic variation and supports planning for both \acp{m-b} and \acp{e-b}.

    \item \textbf{Support for incremental expansion and reconfiguration:} the framework can be applied not only to the design of new \ac{bss} networks but also to the incremental expansion or reconfiguration of existing systems. By explicitly accounting for the existing station layout and its network relationships, the model facilitates planning that strengthens network cohesion.

    \item \textbf{Metaheuristic optimization approach:} a \ac{ga} is used to search for high-quality station configurations while satisfying operational constraints such as the number of stations and the minimum inter-station distance. This approach enables the exploration of diverse scenarios and the generation of near-optimal solutions for large problem instances within reasonable computational time.

\end{enumerate}

To demonstrate the applicability of this framework, we conduct a case study for the city of Barcelona, Spain (see Chapter~\ref{ch:use_case_barcelona} for an overview of the city), using its extensive open data resources. The results illustrate how our approach may help planners systematically explore alternative configurations, assess trade-offs among multiple planning objectives, and design station networks that are tailored to the city's socioeconomic context, infrastructure and topography.

\section{Literature review}

The siting of \ac{bss} stations has been addressed through a wide range of methodological perspectives. Contributions in the literature vary in terms of the data they use, the analytical tools they apply, and the planning objectives they pursue. This section reviews these contributions to clarify the main ways the problem has been addressed and to highlight insights relevant to this study.

    \subsection[From heuristics to digital decision-support tools in BSS planning]{From heuristics to digital decision-support tools in \ac{bss} planning}

Docked \ac{bss} are typically planned through an iterative process, in which goals are first defined and station networks are progressively refined across successive deployment phases~\cite{Duran-Rodas2021, Ferrando2007, Yanocha2018}. Rather than fixing all the station locations at the start, planning guidelines recommend setting clear objectives such as maximizing ridership, promoting social equity, or enhancing first–last kilometer access, and using them to guide preliminary site selection, field validation, and ongoing adjustment. These revisions should be informed by potential demand patterns, operational performance, and stakeholder input~\cite{Ferrando2007, Yanocha2018}. Importantly, such planning goals are not necessarily fixed over time or consistent across a single \ac{bss}. For instance, a \ac{bss} expansion may pursue different priorities than those of the initial implementation.

In practice, early deployments relied on a combination of expert judgment and heuristics. Expert knowledge highlighted key factors shaping cycling demand, such as topography, land use, and connectivity, while heuristics translated these insights into simplified rules that guided site selection. Planning guides reflected this process by recommending station locations in areas of high expected demand, such as near \ac{pt}ation, or commercial hubs, and encouraged the use of \ac{od} surveys and stakeholder consultations to support site selection~\cite{Ferrando2007, Yanocha2018}. These documents emphasized intermodal integration and land-use alignment as core planning principles, often suggesting station placement in visible, accessible, and well-connected locations~\cite{Ferrando2007, Yanocha2018, CROW2006}. Furthermore, bicycle infrastructure manuals underscored the importance of coherent and continuous cycling networks~\cite{AASHTO1999bicycle, CROW2006}. While these guidelines provided valuable practical recommendations, they lacked analytical tools or optimization frameworks capable of systematically evaluating locations or adapting the \ac{bss} to diverse urban conditions. Nonetheless, heuristics and optimization tools can complement each other: practical rules can inform modelling choices, while optimization results can in turn be used to refine or validate heuristic guidelines.

Beyond expert judgment and heuristics, subsequent studies began adopting more systematic methods. As richer spatial data and computational tools became more available in the early 2010s, researchers began to incorporate more formal methods. \ac{gis} were increasingly used to integrate and visualize multiple spatial layers, supporting more data-informed site selection (see Section~\ref{sec:integrating_gis_to_station_placement}). In parallel, spatial optimization models, particularly location–allocation formulations, were introduced to frame station siting as a mathematical problem with explicit objectives and constraints (see Section~\ref{sec:location_allocation_models}). Together, these developments expanded the range of tools available for \ac{bss} planning, combining practical heuristics with data-driven and model-based decision support.


\subsection{Location-allocation models} \label{sec:location_allocation_models}

Location-allocation models are spatial optimization tools that decide where to place a fixed number of service facilities and how to allocate demand to them, with the aim of optimizing one specific performance metric~\cite{Daskin1995}. In these models, demand is typically defined over predefined zones or grid cells, while facilities correspond to candidate sites. In practical applications, these are usually framed as discrete-location problems~\cite{Teixeira2008}, although continuous formulations also exist~\cite{Yeh1996}. These formalizations include the \textit{p}-median problem, which minimizes the total demand-weighted distance to assigned facilities~\cite{Hakimi1965}, and the \ac{mclp}, which seeks to maximize demand served within a defined service radius~\cite{Church1974} (i.e., the maximum distance or travel time from a facility within which users are considered effectively covered). Another influential formulation is the \textit{p}-center problem, which minimizes the maximum distance between any user and their nearest facility~\cite{Hakimi1964}. Depending on the formulation, they can be broadly classified as either efficiency-oriented (e.g., \textit{p}-median), which aim to minimize overall travel effort or system-wide costs, or equity-oriented (e.g., \textit{p}-center),  which aim to ensure fair access by minimizing the longest distance any user must travel to reach a facility~\cite{Murray2010}.

Several studies have applied these classical formulations to \ac{bss} station placement~\cite{Frade2015, Cintrano2018, Lin2011, Mete2018}. In these cases, they typically rely on single-objective optimization frameworks, most often aiming to either maximize spatial coverage or minimize user distance. To ensure a context-sensitive \ac{bss} deployment, models usually incorporate practical constraints such as a fixed number of stations, budget limits, or minimum inter-station distances. For instance, Frade and Ribeiro~\cite{Frade2015} used a variant of the \ac{mclp} to maximize population coverage under budget constraints. Their model optimizes the number, capacity, and location of stations and fleet size, while also considering bike relocations and operational costs. Similarly, Cintrano et al.~\cite{Cintrano2018} studied the \textit{p}-median problem to identify optimal station locations. Lin and Yang~\cite{Lin2011} extended this paradigm by proposing a \ac{minlp} to minimize the total system costs, accounting for user travel distance, station and lane installation, bike inventory, and unmet demand penalties, while ensuring coverage and service availability. Moreover, Mete et al.~\cite{Mete2018} applied several coverage-based formulations (\textit{p}-center, \textit{p}-median, set covering) to locate stations across a university campus, using \ac{gis} to identify and evaluate candidate sites based on proximity to key facilities.

These models have provided reproducible alternatives to heuristic siting but remain focused on classical objectives such as spatial coverage and distance. Broader priorities like social equity or multimodal integration are largely absent, though some aspects have been indirectly considered. For instance, Frade and Ribeiro~\cite{Frade2015} accounted for demand heterogeneity by weighting zones according to estimated trip generation and attraction. This ensured that covering high-demand zones contributed more to the optimization objective than covering low-demand regions, thereby embedding potential usage indirectly rather than treating demand maximization as a distinct objective.


\subsection[Integrating GIS to station placement]{Integrating \ac{gis} to station placement} \label{sec:integrating_gis_to_station_placement}
%Qian2022 té forces referències de \ac{bss} i GIS

Many studies have also incorporated \ac{gis} as a key feature to enrich spatial context and support more informed decision-making~\cite{Fazio2021, Kabak2018, GarciaPalomares2012, Chou2019, Tera2023, Ebrahimi2022}. By integrating and visualizing diverse data layers (e.g., land use, population, slope, and transport networks), \ac{gis} allows planners to evaluate site suitability with greater precision.

One purely \ac{gis}-focused approach is that of Fazio et al.~\cite{Fazio2021}, who proposed a \ac{gis}-based multi-criteria framework to prioritize cycle parking locations in Catania without relying on formal optimization algorithms. Working with a 100×100 meter grid, they integrated spatial data on population, employment, the distribution of \acp{poi}, and \ac{pt} accessibility to calculate thematic indexes for each cell, which were then aggregated into a composite score used to rank and identify the most suitable areas. Similarly, Kabak et al.~\cite{Kabak2018} developed a \ac{gis}-based \ac{mcdm} method combining \ac{ahp} and \ac{moora} to evaluate and rank potential station locations based on twelve spatial criteria, also avoiding formal mathematical optimization but producing explicit suitability maps and priority rankings.

Other studies have combined \ac{gis} with location–allocation models to decide \ac{bss} station locations. García-Palomares et al.~\cite{GarciaPalomares2012} developed one of the earliest \ac{gis}-supported methodologies for siting stations in central Madrid. They estimated potential demand by integrating building-level residential and employment data with zonal trip generation and attraction rates. Using a slope-aware street network to compute bicycle travel times, they applied $p$-median and \ac{mclp} formulations to identify optimal station locations under different network density scenarios.

More recently, Chou et al.~\cite{Chou2019} developed an optimization approach that integrates \ac{gis}-derived spatial data on \ac{pt} networks as key model inputs. Their framework incorporates proportional flow constraints to estimate how bicycles circulate between stations and to assess station capacity requirements. While primarily focused on optimizing fleet allocation and redistribution, their methodology also leverages spatial information to identify candidate station locations and better align supply with latent demand for short-distance trips. Similarly, Tera et al.~\cite{Tera2023} applied a \ac{gis}-based MULTI\ac{moora} model to recommend new station locations in Tartu, Estonia. By integrating spatial data on population density, observed station demand, and proximity to schools and shopping centers, they systematically ranked and compared candidate sites, demonstrating how multi-criteria evaluation can complement network design decisions.

While previous approaches focus on placing new stations, Ebrahimi et al.~\cite{Ebrahimi2022} demonstrated how \ac{gis} can also support the optimization of existing networks. They applied the \ac{mclp} and Target Market Share models to evaluate how well current stations served residential and transit demand, identifying gaps in accessibility and highlighting areas where additional capacity could be prioritized. Although their analysis aimed at optimizing subsets of existing facilities rather than selecting entirely new sites, it illustrates how spatial optimization can inform both operational improvements and broader strategic planning.


\subsection{Multi-criteria and multi-objective station placement}

While many \ac{bss} station placement models have focused on optimizing a single planning goal such as maximizing demand~\cite{GarciaPalomares2012}, maximizing spatial coverage~\cite{Frade2015}, or minimizing user distance~\cite{Cintrano2018}, researchers have increasingly incorporated multiple criteria to better capture the complex factors influencing system performance~\cite{Kabak2018, Guler2021a, Guler2021b, Bahadori2022, Wang2021_modeling, Veillette2018, Fazio2021, Tera2023, Hafezalkotob2019}. In most cases, these \textit{multi-criteria approaches} are used to combine diverse indicators into a single objective function, where all factors are assumed to contribute to the same planning goal. For example, variables such as population density, proximity to \ac{pt}, and cycling infrastructure may be weighted and aggregated to identify sites that best support overall demand. More recently, however, some studies have extended multi-criteria methods to explicitly define multiple planning objectives~\cite{Duran-Rodas2021, Caggiani2020, Fan2024, Qian2022, Nikiforiadis2021}. In these \textit{multi-objective approaches}, different factors are treated as representing distinct goals. This distinction allows planners to explore trade-offs among competing objectives instead of aggregating all criteria into a single composite measure. The following discussion reviews both in detail.

\textit{Multi-criteria approaches} rely on combining diverse spatial and contextual factors into a unified metric. Several strategies have been adopted to determine the importance of each factor, such as the \ac{ahp}, Fuzzy \ac{ahp}, \ac{bwm}, \ac{anp}, or \ac{topsis} to prioritize criteria according to planners' or stakeholders' judgments~\cite{Kabak2018, Guler2021a, Guler2021b, Bahadori2022}. Other studies have used data-driven weighting, where factor importance is derived from statistical models (e.g. regression coefficients or variable importance scores in \ac{ml}) based on observed \ac{bss} usage~\cite{Wang2021_modeling}. A third approach involves flexible weighting, where authors suggest predefined weights to reflect the assumed importance of each factor. For example, Veillette et al.~\cite{Veillette2018} propose a weighting scheme based on existing and potential cyclist demand and proximity to \ac{pt}. However, they explicitly recommend adapting the scheme to match the specific priorities and planning goals of different regions. Similarly, Fazio et al.~\cite{Fazio2021} showcase their methodology using equal weights across all criteria, but emphasise that the weights should ultimately be calibrated in collaboration with local stakeholders and aligned with the intended planning objectives. On the other hand, Tera et al.~\cite{Tera2023} applied the MULTI\ac{moora} model with equal weights across criteria in their evaluation of candidate station locations in Tartu, Estonia. Although the authors did not specify custom weighting in their application, the framework inherently allows planners to adjust weights to reflect specific planning priorities~\cite{Hafezalkotob2019}. 

Multi-criteria approaches typically account for spatial factors such as population and employment density, land-use mix, proximity to \acp{poi}, cycling infrastructure, and access to \ac{pt}. A systematic review by Bahadori et al.~\cite{Bahadori2021} provides an overview of these factors, classifying them into four categories: bike network (e.g., station density, infrastructure), operator (e.g., budget, maintenance costs), user (e.g., walking distance, safety), and city infrastructure (e.g., \ac{pt} connectivity, \acp{poi}, population).

In contrast, \textit{multi-objective approaches} treat factors that pertain to different planning goals separately. A fundamental property of multi-objective optimization problems is that no single solution optimizes all objectives simultaneously; instead, they yield a set of Pareto-optimal solutions that capture the balance between competing objectives. This separation makes the contribution of each objective explicit and enables systematic exploration of alternatives, rather than concealing them within a single aggregated metric~\cite{Rahimi2023, honegumi}. For example, Conrow et al.~\cite{Conrow2018} developed a bi-objective model that maximizes coverage of both residential population and bicycle network infrastructure. 

Several multi-objective studies have incorporated social equity more explicitly. Duran-Rodas et al.~\cite{Duran-Rodas2021} proposed the Demand And/or Equity (DARE) approach, which generates alternative station plans under different emphasis scenarios to illustrate how prioritizing demand, social equity, or a blend of both reshapes spatial allocation. Similarly, Caggiani et al.~\cite{Caggiani2020} integrated multimodal accessibility and social equality into station placement. Their approach minimizes disparities in accessibility among different population groups, measured using a Theil index of multimodal accessibility combining bike-sharing and \ac{pt} travel times. This model balances equality objectives with minimum requirements for system coverage and average accessibility. Moreover, Fan and Harper~\cite{Fan2024} developed a bi-objective integer programming model to jointly maximize coverage of potential demand and improve \ac{pt} for underserved communities. Their approach explicitly weights equity among social groups, showing that even a modest emphasis on equity can yield significant gains in \ac{pt} accessibility for disadvantaged populations with minimal loss in system-wide service quality.

Other multi-objective contributions have placed greater emphasis on explicitly exploring trade-offs along the Pareto frontier. Qian et al.~\cite{Qian2022} focused on the reallocation of existing stations to improve both revenue generation and equitable access to jobs and essential services, particularly for disadvantaged communities. Their framework estimates trip patterns using socioeconomic and built environment data, distributes trips via a gravity model, and employs a \ac{ga} to identify station configurations that balance profitability and social equity objectives within a fixed overall system size. 

Finally, Nikiforiadis et al.~\cite{Nikiforiadis2021} used a tri-objective optimization through a two-phase framework. In the first phase, each objective (maximizing demand coverage, maximizing area coverage, and minimizing rebalancing needs) is optimized independently to establish target benchmark values. In the second phase, these objectives are combined by minimizing the weighted percentage deviations from the targets, enabling planners to adjust the importance of each goal and systematically explore trade-offs. Their approach integrates spatial analysis of \ac{bss} rental patterns and the built environment with practical constraints such as budget, site availability, and minimum station spacing.

Overall, multi-criteria and multi-objective literature on \ac{bss} station placement has substantially advanced the field. Multi-criteria studies have clarified how diverse spatial factors can be assembled into a composite suitability score aligned with a single overarching goal, typically from a predefined set of factors~\cite{Kabak2018, Guler2021a, Guler2021b, Bahadori2022}. Multi-objective formulations build on this by making trade-offs between compiting planning goals explicit and allowing the emphasis across objectives to be adjusted when selecting among the Pareto set~\cite{Duran-Rodas2021, Caggiani2020, Fan2024, Qian2022, Nikiforiadis2021}. Yet in both approaches, sets of spatial factors and planning objectives are commonly fixed \textit{a priori}; adding new dimensions (e.g., safety or environmental criteria) would require reformulating the model. This can constrain the transparent exploration of wider policy scenarios. To adress this, we introduce a modular, adaptable \ac{mcdm} framework that allows new criteria to be incorporated without altering the overall structure. Moreover, planners can flexibly assign weights to these criteria to reflect local priorities, ensuring that different planning objectives can be emphasised as needed.


\subsection{Spatial granularity and network representation}

Studies vary in the spatial scope they consider for planning new \ac{bss} stations, ranging from entire cities to small regions. Many analyses cover a full urban area such as an entire city. For instance, Conrow et al.~\cite{Conrow2018} applied their model to the city of Phoenix and Duran-Rodas et al.~\cite{Duran-Rodas2021} to Munich. In other cases, the focus is on a limited study area. Examples of smaller-scale planning include García-Palomares et al.~\cite{GarciaPalomares2012}, who optimized stations in central Madrid, and Mete et al.~\cite{Mete2018}, who planned stations on a university campus with only 20 candidate sites. 

Equally important is the spatial granularity of potential station locations. Early and some recent optimization models define candidate sites at the zonal level, either by aggregating demand within predefined administrative or planning zones or by applying a uniform grid over the study area. For example, Frade and Ribeiro~\cite{Frade2015} aggregated demand to find potential station locations into predefined neighbourhoods-level regions and Mete et al.~\cite{Mete2018} did similarly using campus sub-areas. Fazio et al.~\cite{Fazio2021} applied a uniform 100×100\,m grid to define candidate cells across the city of Catania, Tera et al.~\cite{Tera2023} and Veillete et al.~\cite{Veillette2018} generated evenly spaced 300\,m grid points covering Tartu and Québec to identify candidate locations, whereas Nikiforiadis et al.~\cite{Nikiforiadis2021} combined grid cells with predefined zones and selected nodes of the transport network. Conrow et al.~\cite{Conrow2018} used census block groups to represent small neighbourhood-scale regions. Duran-Rodas et al.~\cite{Duran-Rodas2021} computed demand and social equity indicators over road network-based catchment zones created by placing regularly spaced virtual station points every 300\,m and generating non-overlapping 300\,m service areas around them using Voronoi tessellation to fully cover the study area. Although these regional approaches simplify the problem and improve tractability, they sacrifice spatial precision, as the exact placement within each cell or zone remains unspecified.

Moving beyond aggregated zones, a second group of studies relies on network-based approaches, selecting discrete candidate points strategically located along major corridors, at intersections of principal streets, or near key \acp{poi} such as \ac{pt} stops. This approach offers more detail than coarse zoning but avoids the computational burden of enumerating every possible location. For instance, Conrow et al.~\cite{Conrow2018}, sampled several hundred potential sites distributed along the bicycle network in their bi-objective model. Similarly, Caggiani et al.~\cite{Caggiani2020} defined their candidate locations as a curated set of around 200 unserved \ac{pt} adjacent nodes that primarily corresponded to bus and metro stops with no existing \ac{bss} stations. Moreover, Fan and Harper~\cite{Fan2024} adopt a similar strategy, generating over 1,200 candidate stations distributed along existing and proposed bike infrastructure, with points placed systematically every 400 metres to ensure coverage of both high-demand and disadvantaged areas. 

To our knowledge, only a few studies pursue high-resolution approaches by using street-network-derived candidate sets for station locations. For example, Cintrano et al.~\cite{Cintrano2018} enumerate an exceptionally large set of 33,550 street segments across Málaga as potential sites, while Xin’s Vancouver analysis~\cite{Xin2013} evaluates 1,489 street intersections within the downtown area. Beyond site definition, these high-resolution approaches and most of the literature rely on network representations in the optimization process to evaluate realistic travel paths rather than simple Euclidean distances or coarse zoning. In these cases, the network is represented as undirected, treating all links as bidirectional with identical distances. Although this simplification does not capture features such as one-way streets and asymmetric accessibility between origins and destinations, these studies show how network-based optimizations can better represent cycling conditions and provide a basis for future refinements.

Additionally, slope is a critical factor in \ac{bss} cycling behaviour, with fewer trips observed on uphill trips and more on downhill ones~\cite{Morency2015, Grau-Escolano2024, MateoBabiano2016}. While Mix et al.~\cite{Mix2022} found that elevation was not a significant predictor in Santiago de Chile, slope has been shown to strongly influence cycling patterns in other cities with heterogeneous topography, such as Barcelona~\cite{Grau-Escolano2024} and Brisbane~\cite{MateoBabiano2016}. In terms of modelling, García-Palomares et al.~\cite{GarciaPalomares2012} took an important step by incorporating slope as a penalty factor uniformly increasing effective travel time. Their approach, however, did not differentiate between uphill and downhill gradients, thereby overlooking the potentially facilitating effect of moderate declines and the safety-related penalties of steep descents.

Our work builds directly on previous high-resolution approaches that derive candidate sites from the street network, such as Cintrano et al.~\cite{Cintrano2018} in Málaga and Xin~\cite{Xin2013} in Vancouver, as well as on the early recognition by García-Palomares et al.~\cite{GarciaPalomares2012} of the need to incorporate slope to achieve a more realistic estimation of user access distances. We extend this line of research by addressing two key limitations. First, unlike prior work that relied on undirected networks, we employ a directed street-level representation that captures circulation constraints such as one-way streets and asymmetric accessibility. This refinement is important because shortest routes in opposite directions can differ substantially in length due to one-way restrictions, with empirical evidence showing asymmetries of more than 10\% in some urban areas~\cite{Melo2022}. Second, we replace the uniform slope penalty of García-Palomares et al. with an effort-based distance measure that differentiates between uphill and downhill gradients. Empirical studies have shown that uphill slopes strongly reduce speed in a near-linear fashion, whereas downhill effects are non-linear: cyclists accelerate only up to a certain gradient and often brake on steeper descents due to safety concerns~\cite{Ryeng2016, Flugel2019}. By incorporating these asymmetries, our approach reflects the increased effort of climbs, the facilitating effect of moderate declines, and the risks associated with steep descents.


\subsection{Spatial arrangement of stations}

In \ac{bss} station planning, a critical design consideration is the spatial arrangement of stations, meaning the overall pattern of how stations are distributed across the service area. Classical network-based formulations such as the $p$-median problem or the $p$-center problem~\cite{Hakimi1964, Hakimi1965} improve accessibility, understood here as reducing the distance between users and their nearest station, by minimizing either average or worst-case user distance. However, in these models the spatial arrangement of facilities remains an emergent outcome rather than an explicit design objective that the planner can control.

The same pattern is evident in \ac{bss}-specific applications of these classical formulations. Although they differ in detail, most focus on improving accessibility: Frade and Ribeiro~\cite{Frade2015} maximized population coverage within a service threshold, Cintrano et al.~\cite{Cintrano2018} minimized average distance, Lin and Yang~\cite{Lin2011} incorporated distance-related costs, and Mete et al.~\cite{Mete2018} explored coverage-based models to balance worst-case distance and facility coverage. While these contributions demonstrate the importance of accessibility in station planning, the spatial arrangement of stations is still treated as a secondary outcome of the optimization rather than a planning objective in its own right.

A complementary dimension of spatial arrangement is proximity, reflecting how close stations are to one another. Accessibility and proximity are distinct dimensions: a tightly clustered set of stations at one edge of the service area may have high inter-station proximity but low accessibility for most users, whereas relocating the same cluster to a central area leaves proximity unchanged while improving accessibility. Considering proximity can be valuable for enabling short trips and facilitating bike rebalancing, yet most models don't explicitly optimize proximity but only impose minimum inter-station distances to avoid excessive clustering of stations~\cite{Nikiforiadis2021, Kabak2018}. In this line, concepts from network science such as farness~\cite{Sabidussi1966}, defined as the sum of shortest-path distances from one node to all others, provide a way to quantify spatial separation among stations within the cycling network. Adapting such measures offers a principled approach to evaluating inter-station spacing, complementing accessibility in capturing the overall structure of the station network.

Building on earlier studies that primarily optimized accessibility~\cite{Frade2015, Cintrano2018, Lin2011, Mete2018} or imposed minimum spacing rules to avoid redundant stations~\cite{Nikiforiadis2021, Kabak2018}, our framework evaluates both accessibility and proximity as complementary dimensions of spatial arrangement. Accessibility is measured by adapting the $p$-median formulation to compute the distance from each node to its nearest station, while proximity is captured by summing pairwise distances among selected stations. These metrics are then combined into a composite score with adjustable weights, making spatial arrangement trade-offs explicit. Moreover, this score works alongside the \ac{mcdm} framework, where additional spatial factors can be added with chosen weights, so that results reflect broader planning goals while balancing accessibility and proximity.


\subsection{Algorithmic approaches}

\ac{gis}-based methods have been used to systematically integrate diverse spatial data~\cite{GarciaPalomares2012, Ebrahimi2022}. Combined with \ac{mcdm} techniques such as \ac{ahp} and \ac{moora}, they allow factors to be weighted and aggregated into composite suitability scores or ranked lists of candidate sites~\cite{Kabak2018, Bahadori2022, Guler2021a, Guler2021b}. Applied on their own, however, \ac{gis} methods mainly support the visualization and ranking of potential sites rather than computing a single optimal solution.

In contrast, optimization models frame station placement as a mathematical problem to maximize or minimize a defined objective under constraints. As described earlier, common formulations include the $p$-median, maximal covering, and $p$-center models. These problems are typically set up as integer programming models with decision variables indicating whether a station is placed at each candidate location~\cite{Frade2015, Lin2011, Mete2018, Cintrano2018}. Constraints often enforce budget limits (or a fixed number of stations), and additional constraints such as minimum spacing between stations and the coverage area of each station are sometimes incorporated in the optimizations~\cite{Conrow2018, Nikiforiadis2021, Caggiani2020}. These models are then solved for an optimal solution using optimization solvers such as CPLEX~\cite{Mete2018, Chou2019}, Gurobi~\cite{Mix2022},  LINGO~\cite{Lin2011}, or FICO Xpress-IVE~\cite{Conrow2018}. While exact solvers guarantee optimality, they can struggle as the candidate pool grows large or the objective function becomes more complex. To address this computational burden, some researchers have used the location-allocation tools in \ac{gis} software such as ArcGIS~\cite{GarciaPalomares2012, Ebrahimi2022}, which implement heuristics for $p$-median or covering problems and can quickly generate near-optimal solutions.

Given the computational burden of exact optimization methods, recent research has increasingly turned to metaheuristic algorithms. Many studies in diverse fields have adopted these methods to search for high-quality solutions within practical computational times~\cite{Tomar2023}. Widely used approaches include \ac{ga}, \ac{simann}, \ac{pso}, and \ac{aco}, all of which have demonstrated strong performance in large-scale combinatorial optimization. These methods have been successfully applied in various domains such as facility location~\cite{Chadry2003, Chalupa2019}, transport planning~\cite{Holliday2023, Sachan2024}, and network design~\cite{Gallo2010, Madadi2024}. Although these metaheuristic approaches were originally designed for single-objective problems, they have often been adapted to multi-objective settings by reformulating the problem into a series of single-objective ones~\cite{Rahimi2023}. Common strategies include weighted aggregation, where objectives are combined into a single fitness function with varying weights, and the $\varepsilon$-constraint method, where one objective is optimized while the others are imposed as constraints with different bounds. Such reformulations, however, usually require multiple runs to approximate the Pareto front and may fail to capture non-convex or discontinuous trade-off surfaces.

To overcome these limitations, dedicated multi-objective metaheuristics have been developed. These algorithms are designed to approximate the Pareto front in a single run by maintaining a diverse set of nondominated solutions, rather than collapsing multiple objectives into one. Well-known examples include \ac{nsgaii}~\cite{Deb2002}, \ac{spea2}~\cite{Zitzler2001} and its extension SPEA2+\cite{Kim2004}, and swarm-based extensions such as \ac{mopso}~\cite{Coello2002} and \ac{monaco}~\cite{Cardoso2010}.

In the context of \ac{bss}, several studies have leveraged metaheuristic methods to optimize station placement. Cintrano et al.~\cite{Cintrano2018} evaluated \ac{simann}, \ac{pso}, Iterated Local Search, and Variable Neighbourhood Search on a $p$-median formulation, finding that the \ac{ga} consistently achieved the best solution quality. Qian et al.~\cite{Qian2022} applied a \ac{ga} in Chicago to balance revenue and social equity objectives, demonstrating its flexibility for planning under competing priorities. Liu et al.~\cite{Liu2015} combined demand prediction using neural networks with a \ac{ga}-based optimization model to select station sites that maximize usage and minimize rebalancing needs in New York City. Caggiani et al.~\cite{Caggiani2020} also implemented a \ac{ga}, highlighting its suitability for equity-focused planning objectives. Finally, Askari et al.~\cite{Askari2017} proposed a hybrid \ac{ga}–\ac{pso} approach to address a capacitated location–allocation problem under demand uncertainty, while Chen et al.~\cite{Chen2020_optimal} developed a hybrid \ac{pso} that incorporates \ac{ga} operators to solve a mixed-integer programming formulation of \ac{e-b} station deployment in Ningbo, China, improving population coverage and accessibility.

Despite the variety of metaheuristic approaches, \acp{ga} have been especially prominent in \ac{bss} research~\cite{Cintrano2018, Liu2015, Caggiani2020, Qian2022}, where their population-based design and flexibility in operator definition make them well suited to the spatial encoding of station locations while also being able to handle multiple planning objectives and constraints. Comparative experiments further suggest that \acp{ga} can outperform some metaheuristic alternatives: Cintrano et al.~\cite{Cintrano2018}, for instance, reported that \ac{ga} produced higher-quality solutions than \ac{simann}, \ac{pso}, Iterated Local Search, and Variable Neighbourhood Search on a $p$-median formulation. On this basis, a \ac{ga} algorithm was adopted in our framework due to: (i) its use in prior \ac{bss} optimization studies; (ii) its ability to incorporate custom, constraint-aware operators; (iii) its suitability for encoding solutions at the street-network level; and (iv) its stronger performance in preliminary tests.

In optimization literature, metaheuristic algorithms are typically validated by evaluating computational time, memory requirements, and the time needed to obtain a reasonably good solution, although the primary focus is often on the solution quality or an error measure~\cite{Osaba2021}. Moreover, because metaheuristics are stochastic, it is standard practice to perform multiple independent runs and report aggregate performance statistics (e.g., best, mean, worst solutions and standard deviations) to assess reliability and variability~\cite{Osaba2021}. The robustness of results across these runs and diverse problem instances is judged by the consistency, ensuring that the algorithm performs well under varying conditions. Furthermore, rigorous comparative studies frequently include statistical significance tests to confirm whether performance differences are meaningful~\cite{Derrac2011, Danka2013}.

Researchers commonly evaluate solution quality by comparing a metaheuristic's best-found objective value against known optima or exact methods when available. Where problem size permits, benchmarking against exact solvers (e.g., \ac{milp} or \ac{ip}) is a standard approach to compute optimality gaps ~\cite{Roshanaei2010}. For instance, in their review of the Travelling Salesman Problem, Alkhalifa et al.~\cite{Alkhalifa2025} describe this gap in terms of how far a solution is from the best possible one, typically reported as a percentage. For larger or more complex instances where exact solutions are impractical, comparisons are conducted against simpler heuristics (e.g., greedy or rule-based methods) or previously published metaheuristic results. In the case where the global optimum is known, several studies report that well-tuned metaheuristics can achieve near-optimal solutions within a few percent of the optimal value~\cite{Helsgaun2000, JohnsonMcGeoch2003, Schwinn2018}.

In contrast to most previous \ac{bss} optimization studies, which rarely benchmark metaheuristic solutions against exact methods or systematically compare them with alternative metaheuristics, our framework evaluates performance by quantifying optimality gaps relative to exact \ac{milp} solutions under simplified objectives. This would provide evidence that the proposed \ac{ga} can deliver competitive solutions in these scenarios.